\chapter{Remaining FYP Plan}

\section{Remaining Scope}

The proposal's nine modules and thirteen objectives remain the scope of the project. At this mid-report four modules are in a working vertical slice (Table~\ref{tab:scope-modules}); the rest, and the evaluation that makes the research claims, remain. Nothing is dropped. The plan below is ordered by what the research questions need first: a measured retrieval baseline, then the evaluation of correction-aware ranking, then the product features around them.

\section{Remaining Modules and Features}

Table~\ref{tab:remaining-modules} states what remains for each proposal module after Mid.

\begin{table}[H]
\centering
\caption{Remaining work by module}
\label{tab:remaining-modules}
\small
\begin{tabularx}{\textwidth}{@{}>{\raggedright\arraybackslash}p{0.22\textwidth} >{\raggedright\arraybackslash}p{0.22\textwidth} L@{}}
\toprule
Module & State now & Remaining \\
\midrule
M01 Workspace & Not started & Authentication, roles, organisations, reviewer identity in the audit trail \\
M02 Assets & Working & Manual asset registration; component view \\
M03 Knowledge graph & Working (one aircraft) & Component-level graph; path exploration; WebGL fallback (PR open) \\
M04 Diagnostics & Working (lexical, heuristic) & Semantic and hybrid retrieval; baselines; server-side session history \\
M05 Conflict intelligence & Working; mostly unanalysed & Token measurement; full-corpus analysis with cascade; review audit log \\
M06 Report intelligence & CSV and single record & Extraction from documents; ambiguous-field flags \\
M07 Academy & Not started & Case-based scenarios, learner steps, feedback linked to evidence \\
M08 Analytics & Not started & Recurring-fault and fleet views with descriptive statistics only \\
M09 Administration & Status page only & Audit, configuration, governance \\
\bottomrule
\end{tabularx}
\end{table}


\section{Remaining Iterations}

The proposal's later iterations stand: Iteration~2 (FYP-1 Final), an integrated minimum viable product with hybrid retrieval, role-based access, multi-turn sessions with server-side history and a deployed demonstration; Iteration~3 (FYP-2 Early) for training, analytics and graph exploration; Iteration~4 (FYP-2 Mid), a job-fair-ready integrated product; and Iteration~5 (FYP-2 Final) for hardening, final evaluation and handover. The percentages in the timeline are targets, not measured progress.

\section{Implementation Plan}

\textbf{Iteration 2 (FYP-1 Final).} Tasks are ordered by dependency; owners follow the ownership in Table~\ref{tab:ownership}.

\begin{table}[!htbp]
\centering
\caption{FYP-1 Final implementation tasks}
\label{tab:fyp1-final-tasks}
\small
\begin{tabularx}{\textwidth}{@{}c L >{\raggedright\arraybackslash}p{0.18\textwidth} >{\raggedright\arraybackslash}p{0.20\textwidth}@{}}
\toprule
\# & Task & Owner & Done when \\
\midrule
1 & Close remaining open defects in Table~\ref{tab:defects} (HTTP 400 for unreadable CSV, orphan record); staging explanations and graph-reset hardening were addressed after \texttt{36a98a3} (Appendix~D) & Usman, Umar & Each has a test or recorded check \\
2 & Measure live token use on a bounded sample and project full-corpus cost; decide default analysis mode & Usman & Measured numbers in an evidence document \\
3 & Embedding-based retrieval and a hybrid ranker behind a switch & Umar & Runs on the corpus; unit tests \\
4 & Diagnostic evaluation harness and baselines (keyword, lexical, semantic, graph-only, hybrid, direct model) & Umar, Usman & Metrics reproducible from a script \\
5 & Roles and authentication; reviewer identity in the audit trail & Usman, Shahmeer & Requests carry an identity; review events record who and when \\
6 & Server-side session history for the Copilot & Umar, Shahmeer & Sessions survive a browser change \\
7 & Analysis of the full corpus with the cascade on a scratch graph, then staging with approval & Usman & Coverage shown in the interface \\
8 & Quality: lint job in CI, coverage report for core modules, release tags & Shahmeer, Usman & CI reports them \\
\bottomrule
\end{tabularx}
\end{table}

\section{Testing and Validation Plan}

\begin{table}[!htbp]
\centering
\caption{Validation plan against the research questions}
\label{tab:validation-plan}
\small
\begin{tabularx}{\textwidth}{@{}c L L@{}}
\toprule
RQ & Question & Planned evidence \\
\midrule
RQ1 & Does graph-constrained retrieval beat keyword retrieval at finding relevant cases? & Recall at $k$ and mean reciprocal rank on the 141 FAA pairs, same queries for both methods \\
RQ2 & Does hybrid retrieval improve ranking over semantic alone? & Same metrics, with an ablation of each signal \\
RQ3 & Does correction-aware retrieval reduce misleading recommendations? & On the 92 of 141 FAA pairs where the follow-up names a different component: does the system surface the later component for the earlier symptom? \\
RQ4 & Can abstention reduce unsupported suggestions? & Coverage against error rate as the evidence threshold changes \\
RQ5 & Does case-based training improve learners' results? & Planned with the training module in FYP-2 \\
\bottomrule
\end{tabularx}
\end{table}

To avoid leakage the evaluation will use a different year of FAA data than the one used to choose weights and thresholds, and any threshold will be fixed before it is applied to held-out data, as the cascade threshold was. Usability will be measured with proxy participants (a task and a System Usability Scale questionnaire), which the proposal already says will be declared as proxies. Load performance will be tested against the 500~ms target (NFR-02), and coverage measured against NFR-04.

\section{Deployment Plan}

The staging deployment, with its recorded rollback and approval gate, is the model for later releases. For FYP-2 the plan is a production-style deployment of the same pipeline with backups of the graph before any migration, read-only demonstrations on shared data, and a prerecorded backup of every demonstration.

\section{Tentative Timeline}

Dates follow the university calendar once it is published; the weeks below are relative to the FYP-1 Final evaluation.

\begin{table}[!htbp]
\centering
\caption{Tentative FYP-1 Final timeline}
\label{tab:timeline}
\small
\begin{tabularx}{\textwidth}{@{}>{\raggedright\arraybackslash}p{0.18\textwidth} L@{}}
\toprule
Period & Work \\
\midrule
Weeks 1--2 & Close open defects (task 1); live token measurement (task 2); lint and coverage in CI (task 8) \\
Weeks 3--5 & Embedding retrieval and hybrid ranker (task 3); evaluation harness and baselines (task 4) \\
Weeks 4--6 & Roles, authentication and reviewer identity (task 5) \\
Weeks 6--7 & Server-side session history (task 6); full-corpus analysis on a scratch graph (task 7) \\
Week 8 & Integration, regression, updated requirements and design specification, rehearsal \\
\bottomrule
\end{tabularx}
\end{table}

Each week ends with a demonstrable milestone, recorded with its evidence in the milestone log, so that progress is checkable against the demonstration and not against a task list.

\section{Expected Final Deliverables}

A deployed, tested and documented multi-user system with the nine modules; a reproducible evaluation against baselines with stated limits; the repository with tags, setup instructions and decision records; the final report; and a demonstration that every member can explain.
